% !TEX root = ../main.tex

\section{Introduction}

What determines the quality of a text? And how do we measure it at scale? These two questions have become increasingly important with the development of advanced Natural Language Generation (NLG) systems  such as LLMs that can readily generate any type of text. Although Generative AI (GenAI) promises great potential for industrial applications, its implementation in business contexts is hindered by the inherent variability of output quality.

One company facing this problem is Barter. Barter is an influencer marketing platform, where companies can post collaboration opportunities ("Deals") which content creators can apply to. GenAI is increasingly being used for marketing copy generation, but without adequate quality control, the accumulation of low-quality output becomes a significant operational problem at high volumes. 

The need for improvements in the validation of NLG systems is increasingly being recognized \citep{xiaoEvaluatingEvaluationMetrics2023, linPromptsConstructsDualValidity2025, binetteImprovingValidityPractical2024, beanMeasuringWhatMatters2025, salaudeenMeasurementMeaningValidityCentered2025, wangEvaluatingGeneralPurposeAI, kardanovaNovelPsychometricsBasedApproach2024, federiakinImprovingLLMLeaderboards2025, PDFModernLLM2026, feuerWhenJudgmentBecomes2025, alaaMedicalLargeLanguage2025}, calling on the methods from psychometrics to formalize the measurement procedure. Psychometrics concerns itself with the measurement of latent, unobservable variables; indeed, in the social sciences, the variable of interest is rarely directly observable. Meaningful measurement therefore requires specifying how the unobservable target construct relates to the indicators or criteria through which it is measured.

Analogously, in text evaluation, complex constructs such as \textit{Text Quality} or \textit{Deal Attractiveness} are not directly readable properties of a text, but can be decomposed into theoretically related subdimensions that enable fine-grained diagnostics. In a seminal paper, \citet{fabbri_summeval_2021} decompose "Summary Quality" into the dimensions \textit{coherence}, \textit{consistency}, \textit{fluency} and \textit{relevance}, which they use to show that summary models have specific failure modes: pre-trained language models correlate highly with human judgments on consistency and fluency, whereas extractive models struggled with coherence and relevance. Multi-dimensional evaluation has become the dominant paradigm in automated evaluation \citep{zhongUnifiedMultiDimensionalEvaluator2022a}, and with the introduction of LLMs, has become increasingly flexible. 

However, many modern evaluation methods do not assume a specific measurement model. As a result, it remains unclear whether the selected dimensions adequately define the construct, how they should be recomposed into an overall score, and what claims that score supports. This highlights the central problem of this thesis: For artefact-level evaluation, when does measurement become meaningful, and under which conditions? 

We argue that artefact-level evaluation can best be understood as a formative measurement problem. In a formative measurement model, the target construct is defined by its constituent dimensions. Because of this, omitting or underspecifying a relevant dimension means omitting or changing part of the construct itself \citep{bollen1991conventional}. This creates a central validity risk: holistic scores, or multi-dimensional scores that are weakly specified, might be ill-aligned with the evaluation target. 


To test this argument, the thesis compares increasingly specified representations of artefact-level quality: holistic scoring, weakly specified or macro-level decomposition, and fine-grained formative decomposition. In the formative condition, theoretically specified dimensions are measured separately and recomposed explicitly into an overall score. This design addresses two research questions:

\begin{description}
\item[\textbf{RQ1:}] Does stronger construct specification provide stronger validity evidence than holistic or weakly specified alternatives?
\item[\textbf{RQ2:}] Does stronger construct specification improve the diagnostic interpretability of artefact-level measurements relative to holistic or weakly specified alternatives?
\end{description}



The analysis is conducted in two validation settings. First, the FeedbackQA dataset is used to assess convergent validity: whether LLM-derived measurements of Q\&A answer quality align with human quality judgments. Second, the Barter Deals dataset is used to assess predictive validity: whether LLM-derived measurements of deal attractiveness predict aggregate creator application behavior. Before comparing validity evidence, the thesis also diagnoses the LLM measurements themselves, examining scale use, rating uncertainty, verbosity sensitivity, and construct separability. Together, these analyses evaluate whether formative decomposition improves artefact-level LLM measurement not only by producing better predictions, but by yielding scores whose interpretation is more explicitly tied to the construct being measured.




